% GENERATED FILE. Edit canonical lessons, then run npm run generate:books.
% canonical-source-hash: fnv1a64:3f0a7b60ce7caa6b
% canonical-lessons: IT-C480-offendersi, IT-C480-rancore, IT-C480-perdono, IT-C480-pentirsi, IT-C480-sfogarsi

\chapter{Perdonare}
\label{ch:it-perdonare}
\hlchaptermodality{480}

\begin{chapteropening}
\textbf{By the end of this chapter:} \emph{I can talk about taking offence, grudges and forgiving.}

Complete the last lesson of chapter 480: I can talk about taking offence, grudges and forgiving.
\end{chapteropening}

\section[offendersi]{offendersi — to take offence}

\label{lesson:IT-C480-offendersi}



\begin{quote}
Before the new one: say the Italian for enthusiasm, then the Italian for to be moved.
\end{quote}

\subsection*{You'll want to know: offendersi}
\textbf{offendersi} — \textquotedblleft{}to take offence\textquotedblright{}.

\begin{quote}
Si è offeso perché non l'abbiamo invitato.
\end{quote}

\begin{quote}
\emph{He took offence because we did not invite him.}
\end{quote}

The past participle is \emph{offeso}. To offend someone: \emph{offendere qualcuno}.

\begin{cousinweb}
\emph{offendersi} is the reflexive of \emph{offendere}, a learned borrowing from Latin \emph{offendere} \textquotedblleft{}to strike against, to displease\textquotedblright{} (Lewis \& Short).
\end{cousinweb}

\begin{grammarlens}[title={What you've built}]
One more word for telling, explaining and describing at B1.
\end{grammarlens}

\subsection*{Your turn}
\begin{itemize}
\raggedright
  \item \emph{Say it:} \emph{offendersi}
  \item \emph{Say it:} \emph{offendersi}, once more
  \item \emph{Say it:} a sentence of your own with \emph{offendersi}
  \item \emph{Recall:} say the Italian for enthusiasm, then the Italian for to be moved, then say \emph{offendersi} again
\end{itemize}

\begin{tcolorbox}[breakable,colback=teal!4,colframe=teal!35!black,title={Before you move on}]
What does offendersi mean? (\textquotedblleft{}To take offence\textquotedblright{}.) Say it once more.
\end{tcolorbox}
\section[il rancore]{il rancore — grudge, resentment}

\label{lesson:IT-C480-rancore}



\begin{quote}
Before the new one: say the Italian for to be moved, then the Italian for to take offence.
\end{quote}

\subsection*{You'll want to know: il rancore}
\textbf{il rancore} — \textquotedblleft{}grudge, resentment\textquotedblright{}.

\begin{quote}
Dopo tanti anni non ho più rancore verso di lui.
\end{quote}

\begin{quote}
\emph{After so many years I no longer bear him a grudge.}
\end{quote}

\emph{Serbare rancore} — to hold a grudge.

\begin{cousinweb}
\emph{rancore} comes from Late Latin \emph{rancōrem} \textquotedblleft{}rancidness, a grudge\textquotedblright{} (EtymDB), from \emph{rancēre} \textquotedblleft{}to be rancid\textquotedblright{} (Lewis \& Short).
\end{cousinweb}

\begin{grammarlens}[title={What you've built}]
One more word for telling, explaining and describing at B1.
\end{grammarlens}

\subsection*{Your turn}
\begin{itemize}
\raggedright
  \item \emph{Say it:} \emph{il rancore}
  \item \emph{Say it:} \emph{il rancore}, once more
  \item \emph{Say it:} a sentence of your own with \emph{il rancore}
  \item \emph{Recall:} say the Italian for to be moved, then the Italian for to take offence, then say \emph{il rancore} again
\end{itemize}

\begin{tcolorbox}[breakable,colback=teal!4,colframe=teal!35!black,title={Before you move on}]
What does il rancore mean? (\textquotedblleft{}Grudge, resentment\textquotedblright{}.) Say it once more.
\end{tcolorbox}
\section[il perdono]{il perdono — forgiveness}

\label{lesson:IT-C480-perdono}



\begin{quote}
Before the new one: say the Italian for to take offence, then the Italian for grudge.
\end{quote}

\subsection*{You'll want to know: il perdono}
\textbf{il perdono} — \textquotedblleft{}forgiveness\textquotedblright{}.

\begin{quote}
Ti chiedo perdono.
\end{quote}

\begin{quote}
\emph{I ask your forgiveness.}
\end{quote}

The noun for \emph{perdonare} \textquotedblleft{}to forgive\textquotedblright{}.

\begin{cousinweb}
\emph{perdono} comes from \emph{perdonare} \textquotedblleft{}to forgive\textquotedblright{}, usually traced to Late Latin \emph{perdōnāre}, built on Latin \emph{dōnāre} \textquotedblleft{}to give\textquotedblright{} (Lewis \& Short).
\end{cousinweb}

\begin{grammarlens}[title={What you've built}]
One more word for telling, explaining and describing at B1.
\end{grammarlens}

\subsection*{Your turn}
\begin{itemize}
\raggedright
  \item \emph{Say it:} \emph{il perdono}
  \item \emph{Say it:} \emph{il perdono}, once more
  \item \emph{Say it:} a sentence of your own with \emph{il perdono}
  \item \emph{Recall:} say the Italian for to take offence, then the Italian for grudge, then say \emph{il perdono} again
\end{itemize}

\begin{tcolorbox}[breakable,colback=teal!4,colframe=teal!35!black,title={Before you move on}]
What does il perdono mean? (\textquotedblleft{}Forgiveness\textquotedblright{}.) Say it once more.
\end{tcolorbox}
\section[pentirsi]{pentirsi — to regret, to repent}

\label{lesson:IT-C480-pentirsi}



\begin{quote}
Before the new one: say the Italian for grudge, then the Italian for forgiveness.
\end{quote}

\subsection*{You'll want to know: pentirsi}
\textbf{pentirsi} — \textquotedblleft{}to regret, to repent\textquotedblright{}.

\begin{quote}
Mi sono pentito di quello che ho detto.
\end{quote}

\begin{quote}
\emph{I regretted what I said.}
\end{quote}

\emph{Pentirsi di} + noun or infinitive.

\begin{cousinweb}
\emph{pentirsi} is inherited from Latin \emph{paenitēre} \textquotedblleft{}to cause regret\textquotedblright{} (EtymDB; Lewis \& Short).
\end{cousinweb}

\begin{grammarlens}[title={What you've built}]
One more word for telling, explaining and describing at B1.
\end{grammarlens}

\subsection*{Your turn}
\begin{itemize}
\raggedright
  \item \emph{Say it:} \emph{pentirsi}
  \item \emph{Say it:} \emph{pentirsi}, once more
  \item \emph{Say it:} a sentence of your own with \emph{pentirsi}
  \item \emph{Recall:} say the Italian for grudge, then the Italian for forgiveness, then say \emph{pentirsi} again
\end{itemize}

\begin{tcolorbox}[breakable,colback=teal!4,colframe=teal!35!black,title={Before you move on}]
What does pentirsi mean? (\textquotedblleft{}To regret, to repent\textquotedblright{}.) Say it once more.
\end{tcolorbox}
\section[sfogarsi]{sfogarsi — to let off steam, to pour one's heart out}

\label{lesson:IT-C480-sfogarsi}



\begin{quote}
Before the new one: say the Italian for forgiveness, then the Italian for to regret.
\end{quote}

\subsection*{You'll want to know: sfogarsi}
\textbf{sfogarsi} — \textquotedblleft{}to let off steam, to pour one's heart out\textquotedblright{}.

\begin{quote}
Quando sono arrabbiato mi sfogo con la corsa.
\end{quote}

\begin{quote}
\emph{When I am angry I let off steam by running.}
\end{quote}

\emph{Lo sfogo} is the outburst; \emph{sfogarsi con qualcuno} — to open up to someone.

\begin{cousinweb}
\emph{sfogarsi} is built from \emph{s-} and \emph{foga} \textquotedblleft{}rush, heat\textquotedblright{}, usually traced to Latin \emph{fuga} \textquotedblleft{}flight, running away\textquotedblright{} (Lewis \& Short).
\end{cousinweb}

\begin{grammarlens}[title={What you've built}]
One more word for telling, explaining and describing at B1.
\end{grammarlens}

\subsection*{Your turn}
\begin{itemize}
\raggedright
  \item \emph{Say it:} \emph{sfogarsi}
  \item \emph{Say it:} \emph{sfogarsi}, once more
  \item \emph{Say it:} a sentence of your own with \emph{sfogarsi}
  \item \emph{Recall:} say the Italian for forgiveness, then the Italian for to regret, then say \emph{sfogarsi} again
\end{itemize}

\begin{tcolorbox}[breakable,colback=teal!4,colframe=teal!35!black,title={Before you move on}]
What does sfogarsi mean? (\textquotedblleft{}To let off steam, to pour one's heart out\textquotedblright{}.) Say it once more.
\end{tcolorbox}
